% Options for packages loaded elsewhere
\PassOptionsToPackage{unicode}{hyperref}
\PassOptionsToPackage{hyphens}{url}
\documentclass[
  ignorenonframetext,
]{beamer}
\newif\ifbibliography
\usepackage{pgfpages}
\setbeamertemplate{caption}[numbered]
\setbeamertemplate{caption label separator}{: }
\setbeamercolor{caption name}{fg=normal text.fg}
\beamertemplatenavigationsymbolsempty
% remove section numbering
\setbeamertemplate{part page}{
  \centering
  \begin{beamercolorbox}[sep=16pt,center]{part title}
    \usebeamerfont{part title}\insertpart\par
  \end{beamercolorbox}
}
\setbeamertemplate{section page}{
  \centering
  \begin{beamercolorbox}[sep=12pt,center]{section title}
    \usebeamerfont{section title}\insertsection\par
  \end{beamercolorbox}
}
\setbeamertemplate{subsection page}{
  \centering
  \begin{beamercolorbox}[sep=8pt,center]{subsection title}
    \usebeamerfont{subsection title}\insertsubsection\par
  \end{beamercolorbox}
}
% Prevent slide breaks in the middle of a paragraph
\widowpenalties 1 10000
\raggedbottom
\AtBeginPart{
  \frame{\partpage}
}
\AtBeginSection{
  \ifbibliography
  \else
    \frame{\sectionpage}
  \fi
}
\AtBeginSubsection{
  \frame{\subsectionpage}
}
\usepackage{iftex}
\ifPDFTeX
  \usepackage[T1]{fontenc}
  \usepackage[utf8]{inputenc}
  \usepackage{textcomp} % provide euro and other symbols
\else % if luatex or xetex
  \usepackage{unicode-math} % this also loads fontspec
  \defaultfontfeatures{Scale=MatchLowercase}
  \defaultfontfeatures[\rmfamily]{Ligatures=TeX,Scale=1}
\fi
\usepackage{lmodern}
\ifPDFTeX\else
  % xetex/luatex font selection
\fi
% Use upquote if available, for straight quotes in verbatim environments
\IfFileExists{upquote.sty}{\usepackage{upquote}}{}
\IfFileExists{microtype.sty}{% use microtype if available
  \usepackage[]{microtype}
  \UseMicrotypeSet[protrusion]{basicmath} % disable protrusion for tt fonts
}{}
\makeatletter
\@ifundefined{KOMAClassName}{% if non-KOMA class
  \IfFileExists{parskip.sty}{%
    \usepackage{parskip}
  }{% else
    \setlength{\parindent}{0pt}
    \setlength{\parskip}{6pt plus 2pt minus 1pt}}
}{% if KOMA class
  \KOMAoptions{parskip=half}}
\makeatother
\usepackage{graphicx}
\makeatletter
\newsavebox\pandoc@box
\newcommand*\pandocbounded[1]{% scales image to fit in text height/width
  \sbox\pandoc@box{#1}%
  \Gscale@div\@tempa{\textheight}{\dimexpr\ht\pandoc@box+\dp\pandoc@box\relax}%
  \Gscale@div\@tempb{\linewidth}{\wd\pandoc@box}%
  \ifdim\@tempb\p@<\@tempa\p@\let\@tempa\@tempb\fi% select the smaller of both
  \ifdim\@tempa\p@<\p@\scalebox{\@tempa}{\usebox\pandoc@box}%
  \else\usebox{\pandoc@box}%
  \fi%
}
% Set default figure placement to htbp
\def\fps@figure{htbp}
\makeatother
\setlength{\emergencystretch}{3em} % prevent overfull lines
\providecommand{\tightlist}{%
  \setlength{\itemsep}{0pt}\setlength{\parskip}{0pt}}
% Auto-generated by infrastructure/rendering/_pdf_latex_helpers.py::extract_math_font_preamble.
% Minimal header passed to Pandoc via -H so Beamer slides pick up the same math font
% as the combined-PDF path. Adjust the manuscript's preamble.md to change the font globally.
\usepackage{unicode-math}
\setmathfont{latinmodern-math.otf}
\usepackage{bookmark}
\IfFileExists{xurl.sty}{\usepackage{xurl}}{} % add URL line breaks if available
\urlstyle{same}
\hypersetup{
  hidelinks,
  pdfcreator={LaTeX via pandoc}}

\author{\texorpdfstring{}{}}
\date{}

\begin{document}

\section{Active Inference as a Process Theory of
Case}\label{sec:cognitive-integration}

\begin{frame}[allowframebreaks]{Active Inference as a Process Theory of
Case}
\textbf{Where we are in the argument.}
\autoref{sec:case-systems}--\autoref{sec:topos-theory} have given the
framework a \emph{structural} theory --- objects, morphisms, functors,
enriched weights, and topos-theoretic bridges. This chapter supplies the
missing \emph{process} theory: active inference recasts case assignment
as variational Bayesian inference over a generative model whose state
variable is the case-role posterior, whose precision parameters are the
enriched weights of \autoref{sec:enriched-categories}, and whose
free-energy descent turns ``parsing a sentence'' into a sequence of
belief updates --- the dynamics that
\autoref{sec:diagrammatic-cognition} exploits to derive ERP predictions.
\end{frame}

\begin{frame}[allowframebreaks]{Static Categories Are Not Enough}
\protect\phantomsection\label{static-categories-are-not-enough}
The preceding sections constructed a mathematical infrastructure for
analyzing case systems---categorical, type-logical, distributional,
enriched, and topos-theoretic. Yet these frameworks remain
\emph{static}: they describe the structure of case grammar without
explaining how a cognitive agent \emph{deploys} that structure during
real-time comprehension and production. Bridging this gap requires a
dynamic \emph{process theory} of case-marked relational reasoning.
Active inference {[}@namjoshi2026fundamentals; @friston2017active{]}
provides exactly this missing dynamic computational layer.
\end{frame}

\begin{frame}[fragile,allowframebreaks]{Surprise Minimization Drives
Case-Frame Inference}
\protect\phantomsection\label{surprise-minimization-drives-case-frame-inference}
\begin{block}{Free Energy Bounds Surprisal}
\protect\phantomsection\label{free-energy-bounds-surprisal}
Active inference is the primary process theory derived from the free
energy principle (FEP): every self-organizing system maintains its
structural integrity by minimizing the surprisal (negative
log-probability) of its sensory observations under an internal
\emph{generative model} of its environment {[}-@friston2010free{]}. The
system executes this minimization through two complementary strategies:

\begin{enumerate}
\tightlist
\item
  \textbf{Perceptual inference}: Update internal beliefs to better
  predict current observations (reduce prediction error)
\item
  \textbf{Active inference}: Act on the environment to bring
  observations in line with predictions (reduce expected prediction
  error)
\end{enumerate}

Recent extensions of active inference to linguistics and cognitive
science have modeled language comprehension and production as forms of
sequential Bayesian inference. As Donnarumma, Frosolone, and Pezzulo
(2023) demonstrate in their integration of large language models and
active inference for modelling eye movements in reading, linguistic
processing constitutes ``inference over a hierarchical generative model,
facilitating predictions and inferences at various levels of
granularity, from syllables to sentences''
{[}-@donnarumma2023integrating{]}. Similarly, Friston et al.~(2021) have
demonstrated how communication emerges between synthetic subjects:
``linguistic outcomes (specifically, the spoken word)\ldots{} are
selected to minimise the free energy given current beliefs'' via
``high-order interactions among abstract (discrete) states in deep
(hierarchical) models'' {[}-@friston2021understanding;
-@friston2020generative{]}.

Both strategies minimize the same mathematical quantity---variational
free energy---and both draw from a single generative model encoding the
system's prior expectations about the relational structure of its world.

Critically, recent neurolinguistic evidence directly supports this
prediction-error account. Li and Futrell {[}-@li2023decomposition;
-@li2024shallow{]} decompose surprisal into two orthogonal components:
\emph{heuristic surprise} (``shallow surprisal''), which tracks the N400
brain potential and reflects lexical-associative prediction error, and a
\emph{discrepancy signal} (``deep surprisal''), which tracks the P600
and reflects structural reanalysis when the true parse diverges from the
initially inferred structure. This decomposition maps directly onto our
enriched case framework: the N400 corresponds to distributional
prediction error \emph{within} a case-role subspace (semantic mismatch),
while the P600 corresponds to structural prediction error \emph{between}
case-diagram topologies (morphosyntactic reanalysis requiring a change
in the generative model's case assignments). The formal equations in
\autoref{sec:daif-results}
(\autoref{eq:eq-7c-n400}--\autoref{eq:eq-7c-p600}) instantiate exactly
this dual decomposition.

\begin{block}{Generative Models of Relational Structure}
\protect\phantomsection\label{generative-models-of-relational-structure}
Under this paradigm, language understanding manifests as \emph{active
inference over relational structure}: the listener maintains a
generative model anticipating who-does-what-to-whom, and each incoming
word supplies evidence that updates this model. Morphological case
marking provides high-precision evidence---for example, a nominative
suffix predicts that the noun phrase functions as the agent, reducing
uncertainty about the relational structure of the unfolding event.
\end{block}

\begin{block}{S-HAI: The Case Diagram as the Abstract ``Schema'' Level}
\protect\phantomsection\label{s-hai-the-case-diagram-as-the-abstract-schema-level}
These relational generative models find their formal articulation in
recent advances such as \textbf{Schema-Based Hierarchical Active
Inference (S-HAI)} {[}@maele2026schema{]}. Unifying predictive
processing with schema theory, S-HAI employs a dual-level POMDP
structure to model rapid generalization across environments. In the
linguistic domain, the ``Level 2'' model encodes abstract, hidden
relational goals---which corresponds exactly to the \emph{case diagram
structure} we describe here. The ``Level 1'' model encodes concrete
sensorimotor navigation---for linguistics, this maps to the sequential
parsing of surface word forms.

Just as S-HAI explains sudden ``zero-shot'' behavioral remapping in
novel environments by preserving the high-level schema mapping while
updating the ``grounding likelihoods'' to new observables, a case frame
enables an agent to rapidly generalize the relational structure of a
complex sentence regardless of novel vocabulary pairings. The abstract
string diagram is the schema; case inflection is the grounding
likelihood.
\end{block}
\end{block}

\begin{block}{The Five-Step
Prior--Observation--Update--Prediction--Action Loop}
\protect\phantomsection\label{the-five-step-priorobservationupdatepredictionaction-loop}
The process unfolds as follows:

\begin{enumerate}
\tightlist
\item
  \textbf{Prior}: The listener has a prior belief about the relational
  structure (a ``case diagram'' encoding expected roles and their
  connections)
\item
  \textbf{Observation}: Each word provides sensory evidence---its form,
  its case marking, its distributional properties
\item
  \textbf{Update}: The listener updates the case diagram to accommodate
  the evidence, using approximate Bayesian inference (typically
  variational message passing)
\item
  \textbf{Prediction}: The updated diagram generates predictions about
  upcoming words (case-marked NPs, verb valency patterns)
\item
  \textbf{Action}: In production, the speaker selects words and case
  markers that minimize expected free energy---choosing expressions that
  are informative, contextually appropriate, and syntactically
  well-formed
\end{enumerate}
\end{block}

\begin{block}{Case Diagrams as Instantiated Situations}
\protect\phantomsection\label{case-diagrams-as-instantiated-situations}
This dynamic active inference perspective connects naturally to
\textbf{situation semantics} (Barwise and Perry
{[}-@barwise1983situations{]}), which treats linguistic meaning as
structured situations---specific configurations of individuals,
relations typed by arity, and spatiotemporal locations, grounded in an
ecological realism where meanings are recurring relational patterns that
organisms attune to. Translated into our categorical framework, a
\textbf{situation} is an instantiated case diagram: a specific
assignment of entities to roles with particular morphisms activated; the
\textbf{situation type} is the structural case category itself, the
abstract pattern that any particular situation instantiates; and
\textbf{information flow} between situations is a functorial mapping
between case categories. Where classical situation semantics left the
dynamics implicit, active inference supplies the computational engine:
the agent moves through situations in real time, updating its case
diagram with each incoming word and using the updated diagram to predict
which situation will arise next.
\end{block}

\begin{block}{Belief Dynamics Over Competing Case Frames}
\protect\phantomsection\label{belief-dynamics-over-competing-case-frames}
\autoref{fig:active-inference-belief} shows a minimal
\textbf{scalar-belief} simulation: the agent holds a
\texttt{CaseDiagramBelief} over alternative alignment frames (NOM--ACC
vs.~ERG--ABS). As syntactic evidence arrives, variational free energy
and entropy track the discrete update loop that
\autoref{sec:daif-results} extends to full return distributions in DAIF.
Generated programmatically from
\texttt{src/visualization/active\_inference\_plots.plot\_alignment\_frame\_belief\_dynamics()}
(belief trajectory from \texttt{sequential\_belief\_update()} in
\texttt{src/cognitive/belief\_updating.py}).

\begin{figure}
\centering
\pandocbounded{\includegraphics[keepaspectratio,alt={Variational free energy drives convergence to the correct case frame during belief updating. The agent begins with a uniform prior over possible case frames (NOM--ACC vs.~ERG--ABS, H{[}q{]} = \textbackslash log 2). Top: stacked P(\textbackslash text\{frame\}) over evidence steps (including the prior column). Middle: entropy H{[}q{]} in nats, with the steepest drop annotated at the most informative step. Bottom: variational free energy F{[}q{]} after each update (per-step curve) and its running minimum \textbackslash min\_\{\textbackslash tau\textbackslash leq t\} F{[}q\_\textbackslash tau{]} (non-increasing envelope), with dashed vertical markers at each evidence index (same convention as word arrivals in ). Likelihoods are synthetic categorical draws consistent with a TQNN-evaluated diagram. Generated programmatically from src/visualization/active\_inference\_plots.plot\_alignment\_frame\_belief\_dynamics() (PNG from scripts/generate\_cognitive\_figures.py).}]{output/figures/active_inference_belief.png}}
\caption{Variational free energy drives convergence to the correct case
frame during belief updating. The agent begins with a uniform prior over
possible case frames (NOM--ACC vs.~ERG--ABS, \(H[q] = \log 2\)).
\textbf{Top}: stacked \(P(\text{frame})\) over evidence steps (including
the prior column). \textbf{Middle}: entropy \(H[q]\) in nats, with the
steepest drop annotated at the most informative step. \textbf{Bottom}:
variational free energy \(F[q]\) after each update (per-step curve) and
its running minimum \(\min_{\tau\leq t} F[q_\tau]\) (non-increasing
envelope), with dashed vertical markers at each evidence index (same
convention as word arrivals in \autoref{fig:daif-free-energy}).
Likelihoods are synthetic categorical draws consistent with a
TQNN-evaluated diagram. Generated programmatically from
\protect\texttt{src/visualization/active\_inference\_plots.plot\_alignment\_frame\_belief\_dynamics()}
(PNG from
\protect\texttt{scripts/generate\_cognitive\_figures.py}).}\label{fig:active-inference-belief}
\end{figure}
\end{block}
\end{frame}

\end{document}
